\subsection{Knowledge-aware Recommendation}
Knowledge-aware recommendation has evolved significantly, aiming to alleviate data sparsity by incorporating external KGs. Early embedding-based methods~\cite{zhang2016collaborative,cao2019unifying,wang2018dkn} primarily utilized TransR or auto-encoders to extract semantic features from KGs to enrich item representations.
Subsequently, propagation-based methods~\cite{wang2019kgat,wang2021learning,yang2022knowledge,yang2023knowledge,zou2022improving,zou2024knowledge,li2025hypercomplex,li2025lightkg} emerged as the dominant paradigm. For instance, KGAT~\cite{wang2019kgat} employs graph attention mechanisms to recursively propagate embeddings across high-order neighbors, explicitly modeling connectivity. KGIN~\cite{wang2021learning} further refines this by disentangling user intents behind interactions via relational paths.
To tackle the inherent noise in KGs, recent advancements have introduced CL~\cite{zou2022improving,yang2023knowledge}. KGIC~\cite{zou2022improving} and KGRec~\cite{yang2023knowledge} construct multi-view contrastive tasks or leverage rationalization scores to denoise supervision signals.
However, most existing approaches adopt a cascaded architecture where KG embeddings are aggregated before collaborative interaction modeling. This tight coupling inevitably propagates structural noise (e.g., high-entropy relations) into the user preference learning process, limiting the robustness of the final representations.

\subsection{Generative Models for Recommendation}
Generative models, particularly diffusion models, have recently garnered attention in recommender systems for their superior ability to model complex distributions~\cite{lin2025diffusion,wei2025diffusion}. Methods like DiffKG~\cite{jiang2024diffkg} and KMDCL~\cite{li2025mask} apply diffusion processes to reconstruct KGs or generate denoised views, achieving state-of-the-art performance. Despite their success, diffusion models suffer from slow inference speeds due to multi-step SDE sampling and often require computationally expensive graph-level generation.
To overcome these efficiency bottlenecks, FM~\cite{lipman2023flow} has been introduced as a more efficient Continuous Normalizing Flow paradigm based on deterministic ODEs. Recent works have successfully adapted FM to various recommendation scenarios: FlowCF~\cite{liu2025flow} applies discrete flow matching to collaborative filtering; FlowRec~\cite{li2025preference} and FMREC~\cite{ijcai2025p346} utilize FM to model dynamic user preference trajectories in sequential recommendation; RecFlow~\cite{huangflow} addresses anisotropic noise in social recommendation.
Nevertheless, the potential of Flow Matching in Knowledge-aware Recommendation remains unexplored. Existing FM-based recommenders typically focus on homogeneous graphs or user sequences, failing to address the unique challenges of KGs, such as the heterogeneity of relations and the prevalence of one-to-many structural noise.